\section{Finite-difference penalized trend estimation}

\subsection{Penalized least squares}

At forecast origin $T$, let
\[
x_T=(y_{T-L+1},\ldots,y_T)^\top\in\mathbb{R}^{L}
\]
denote the estimation window. For difference order $d<L$, let
\[
D_d\in\mathbb{R}^{(L-d)\times L},
\qquad
Q=D_d^\top D_d.
\]
The baseline estimator solves
\begin{equation}
\label{eq:pls_problem}
\widehat\tau_T(\lambda)
=
\arg\min_{\tau\in\mathbb{R}^{L}}
\left\{
\|x_T-\tau\|_2^2+\lambda\|D_d\tau\|_2^2
\right\},
\qquad \lambda\ge0,
\end{equation}
with solution
\begin{equation}
\label{eq:H}
\widehat\tau_T(\lambda)
=
H_\lambda x_T,
\qquad
H_\lambda=(I+\lambda Q)^{-1}.
\end{equation}

The paper uses the zero-drift, identity-weight formulation in \eqref{eq:pls_problem} to isolate the smoothness-selection question. Guerrero's broader statistical formulation allows nonzero mean differences and more general covariance weighting \cite{guerrero2007penalized}; these are natural extensions but are not silently imposed in the baseline criterion.

\subsection{Spectral representation}

Because
\[
v^\top Qv=\|D_dv\|_2^2\ge0,
\]
$Q$ is symmetric positive semidefinite. Under the standard difference operator,
\[
\operatorname{rank}(Q)=L-d,
\]
so $Q$ has $d$ zero eigenvalues and $L-d$ positive eigenvalues. Write
\[
Q=
U\operatorname{diag}
\left(
\underbrace{0,\ldots,0}_{d},
\delta_1,\ldots,\delta_{L-d}
\right)U^\top,
\qquad
\delta_j>0.
\]
Then
\begin{equation}
\label{eq:H_spectral}
H_\lambda
=
U\operatorname{diag}
\left(
\underbrace{1,\ldots,1}_{d},
\frac{1}{1+\lambda\delta_1},
\ldots,
\frac{1}{1+\lambda\delta_{L-d}}
\right)U^\top.
\end{equation}

Thus the $d$ null-space directions are reproduced without shrinkage, while each penalizable direction is attenuated by
\[
\alpha_j(\lambda)=\frac{1}{1+\lambda\delta_j}.
\]
For $d=2$, the null space consists of discrete linear sequences, so increasing $\lambda$ moves the fitted path toward a linear trend.

\subsection{Normalized smoothness}

Guerrero's trace-based construction motivates measuring smoothing through the effective action of $H_\lambda$. We use
\begin{equation}
\label{eq:S}
S(\lambda)
=
1-
\frac{1}{L-d}
\sum_{j=1}^{L-d}
\frac{1}{1+\lambda\delta_j}.
\end{equation}

The normalization excludes the $d$ directions that cannot be penalized, so
\[
S(0)=0,
\qquad
\lim_{\lambda\to\infty}S(\lambda)=1.
\]

For a fixed linear smoother,
\[
\operatorname{edf}(\lambda)=\operatorname{tr}(H_\lambda).
\]
Combining \eqref{eq:H_spectral} and \eqref{eq:S},
\begin{equation}
\label{eq:edf}
\boxed{
\operatorname{edf}(\lambda)
=
L-(L-d)S(\lambda).
}
\end{equation}
Hence $S$ is a normalized complement of effective degrees of freedom.
